% ch06a —— 原书 p.233–240（章首 + §6.1、§6.2 标题）
% [原文 p.233]
\chapter{三维翻模}

本章开始对三维双有理几何作详细研究。与代数曲面相比，三维几何的主要特征是翻转与翻模的出现。本章专门研究翻模。

§6.1 是翻转与翻模的一般入门。这些结果在所有维数下都成立。

具有终端奇点的三维簇的翻模，在 §6.2 中处理。翻模的存在性由指数 1 的三维终端奇点的分类证明。这是四维以上的翻模的存在性尚不清楚的原因之一。翻模链也一并处理。从应用的立场看，确立改良簇的 MMP 过程不会陷入翻转或翻模的无限链，是重要的事情。

§6.3 是此前三维标准奇点研究的继续。主定理 6.23 与定理 6.25 表明：如何通过反复进行具体的、可以看清的爆破，把标准奇点单纯化，直到得到只具有 Q-分解终端奇点的簇。其证明基于 §5.3 中证明的指数 1 的标准奇点的分类。这些 Q-分解终端化的存在性，在假定三维 MMP 时即可得到。然而，这些终端化是确立 MMP 的一部分所必需的，所以不能假定三维 MMP 再加以证明。

§6.4 中，用 Q-分解终端化构成具有标准奇点的三维簇的翻模。
% [原文 p.234]
称为 crepant 降下（crepant descent）的这一方法，把标准翻模的存在性归约为终端翻模。这一方法在许多不同的状况以及更高维下也能起作用。

\section{翻转与翻模}
\takeaway{本节说明翻转与翻模的一般结果，主要是唯一性，以及翻转与某类代数的有限生成性之间的关系；至于翻转与翻模的存在性，即以后各节的主题，所知甚少。}

本节说明关于翻转与翻模的一般结果。主要处理的是唯一性，以及翻转与某类代数的有限生成性之间的关系。关于后继各节的主题即翻转与翻模的存在性，所知甚少。

\begin{remark}{6.1}
本节中，所有的定义与定理等，都是对正规簇叙述的。但是，在下面列举的设定下，定义仍有意义、结果也成立，这一点应当认识到。

（1）有限型的概形（包含混合特征的情形）；

（2）复解析空间。

后者的情形下，把讨论限制在紧集的邻域内是方便的。
\end{remark}

% [Tutor Note: 本章在代数簇与复解析空间两种设定下并行展开；解析情形须把讨论限制在紧集邻域内，故“紧集”“邻域”等词保留原意，不简化为“开集”。]

$(K+D)$-翻转的一般定义已在定义 3.33 中给出。

[Kaw88] 中的下述引理，把翻转与相伴代数的有限生成性联系起来，并且还意味着翻转的唯一性。唯一性的其他证明方法则请参看 [Kol91b, 2.1]。

\begin{lemmap}{6.2}
设 $Y$ 为正规代数簇（或解析空间），$B$ 为 $Y$ 上的 Weil 除子。此时下述两条等价。

\begin{enumerate}[label=（\arabic*）,leftmargin=2.8em]
\item $R(Y,B):=\bigoplus_{m\ge0}\cO_Y(mB)$ 是有限生成 $\cO_Y$-代数层。
\item 适当选取射影的（或固有的）双有理态射 $g:Z\to Y$，可使 $Z$ 正规、$\operatorname{Ex}(g)$ 的余维数至少为 2，且 $B':=g_*^{-1}B$ 是 Q-Cartier 的，并在 $Y$（或 $Y$ 的任意紧子集）上 $g$-丰富。
\end{enumerate}

具有上述性质的 $g:Z\to Y$，除显然的同构外只有一个。
\end{lemmap}

% [原文 p.235]
\begin{proof}
解析的情形下，在 $Y$ 的任意紧子集的邻域内讨论。

假定 (2)。首先证明 $\cO_Y(mB)=g_*\cO_Z(mB')$。显然有自然的单射 $g_*\cO_Z(mB')\to\cO_Y(mB)$。设 $U\subset Y$ 为开集，任取 $\sigma\in H^0(U,\cO_Y(mB))$。$\sigma$ 可延拓为 $\sigma'\in H^0(g^{-1}(U)\setminus\operatorname{Ex}(g),\cO_Z(mB'))$。$\operatorname{Ex}(g)$ 的余维数至少为 2，故 $\sigma'$ 可延拓为截面 $g^*\sigma\in H^0(g^{-1}(U),\cO_Z(mB'))$。因此 $g_*\cO_Z(mB')\to\cO_Y(mB)$ 也是满射。由于 $B'$ 是 $g$-丰富的，$\bigoplus_{m\ge0}\cO_Y(mB)=\bigoplus_{m\ge0}g_*\cO_Z(mB')$ 是有限生成的。

特别地，$Z=\Proj_Y\bigoplus_{m\ge0}\cO_Y(mB)$，故 $Z$ 只有一个。

假定 (1)。以 $B$ 的倍取代换 $B$，可假定 $\cO_Y(B)$ 生成 $\cO_Y$-代数 $R(Y,B)$。令 $Z=\Proj_Y R(Y,B)$。此时 $\cO_Z(1)$ 十分丰富，且 $g_*\cO_Z(m)=\cO_Y(mB)$ 成立。假定 $g:Z\to Y$ 具有例外除子 $E$。此时 $\cO_Z\subseteq\cO_Z(E)$，且存在单射
\[
\tau:\cO_Y(mB)=f_*\cO_Z(m)\subseteq f_*(\cO_Z(m)(E))\qquad(m\gg0).
\]
然而 $\cO_Y(mB)$ 是自反的，且 $\tau$ 在余维数至少为 2 的子集 $g(\operatorname{Ex}(g))$ 之外是同构，这是矛盾。
\end{proof}

% [Tutor Note: 上式中 $\tau$ 的右侧，页图印作 $f_*$，但按文意应为 $g_*$，英文底本亦作 $g_*$；日文本此处疑为笔误，正文照抄不改，仅加此注。]

\begin{remark}{6.3}
设 $Y=\Spec S$ 为仿射簇，$B$ 是由理想 $I\subset S$ 定义的有效除子。此时 $R(Y,-B)=\bigoplus_{m\ge0}I^{(m)}$ 成立。其中 $I^{(m)}$ 表示 $I$ 的 $m$ 次符号幂。$R(Y,-B)$ 也称为 $I$ 的符号幂代数（symbolic power algebra）。
\end{remark}

这样的代数成为有限生成的情形，已知的非常少。[Ree58] 与 [Cut88b] 给出了不成为有限生成的例子。

\begin{corollaryp}{6.4}
设 $f:X\to Y$ 为 $(K+D)$-翻转收缩。即 $X,Y$ 为正规簇，$f:X\to Y$ 为 $\operatorname{Ex}(f)$ 在 $X$ 内的余维数至少为 2 的固有双有理态射，而 $D$ 为 $X$ 上的 Q-除子，且 $-K_X-D$ 为 Q-Cartier 且 $f$-丰富。选取 $r\in\Z_{>0}$，使 $rD$ 成为（整系数）除子。此时下述成立。

\begin{enumerate}[label=（\arabic*）,leftmargin=2.8em]
\item $f$ 的 $(K+D)$-翻转，与 $Y$ 上 $(X,D)$ 的拟标准模型（定义 3.50）一致。
% [原文 p.236]
\item $f$ 的 $(K+D)$-翻转存在 $\iff$ $\bigoplus_{m\ge0}\cO_Y(mf_*(rK_X+rD))$ 是有限生成 $\cO_Y$-代数层。
\item $f$ 的 $(K+D)$-翻转只有一个。
\item 对于使 $K_X+D'$ 成为 Q-Cartier 的 $X$ 上的 Q-除子 $D'$，假定存在 $a,a'\in\Z_{>0}$ 使 $a(K_X+D)\sim_f a'(K_X+D')$。此时 $(K+D)$-翻转也是 $(K+D')$-翻转。
\end{enumerate}
\end{corollaryp}

\begin{proof}
比较定义即得 (1)，再由定理 3.52 得 (3)。唯一性也可由令 $B=f_*(K_X+D)$ 后应用引理 6.2 得到。

(2) 也可由定理 3.52 或引理 6.2(2) 得到。

若 $a(K_X+D)\sim_f a'(K_X+D')$，则 $a(K_{X^+}+D^+)\sim_{f^+}a'(K_{X^+}+D'^+)$，且 $K_{X^+}+D'^+$ 也是 $f^+$-丰富的。
\end{proof}

% [Tutor Note: 推论 6.4 是本节的主干：它把翻转的存在性化为某个分次代数的有限生成性，从而把几何问题变成代数问题；条件 (1) 中的“拟标准模型”见定义 3.50。]

\begin{definition}{6.5}
设 $f:X\to Y$ 为射线型翻转收缩（定义 3.34）。此时由推论 6.4(3)，$(K+D)$-翻转不依赖于 $D$ 的选取。这种情况下，$f^+:X^+\to Y$ 称为 $f$ 的翻转（flip）。
\end{definition}

\begin{proposition*}{6.6}
设 $Y$ 为正规簇，$B$ 为 $Y$ 上的 Weil 除子。对闭点 $0\in Y$，以 $\hat\cO_{0,Y}$ 表示局部环 $\cO_{0,Y}$ 关于极大理想的完备化，并令 $\hat Y=\Spec\hat\cO_{0,Y}$。此时下述两条等价。

\begin{enumerate}[label=（\arabic*）,leftmargin=2.8em]
\item $\bigoplus_{m\ge0}\hat\cO_{0,Y}(mB|_{\hat Y})$ 是有限生成 $\hat\cO_{0,Y}$-代数。
\item 存在 $0$ 的 Zariski 开邻域 $U\subset Y$，使 $\bigoplus_{m\ge0}\cO_U(mB|_U)$ 成为有限生成 $\cO_U$-代数。
\end{enumerate}
\end{proposition*}

\par\medskip\noindent\textbf{注意}\quad 已知 $\hat\cO_{0,Y}$ 也是正规的，但本书不需要这一点。

\begin{proof}
首先，(2) 推出 (1) 可由下述同构得到。
\[
\bigoplus_{m\ge0}\hat\cO_{0,Y}(mB|_{\hat Y})=\hat\cO_{0,Y}\otimes\Bigl(\bigoplus_{m\ge0}\cO_U(mB|_U)\Bigr).
\]
% [原文 p.237]
反之，假定 $\bigoplus_{m\ge0}\hat\cO_{0,Y}(mB|_{\hat Y})$ 由次数不超过 $n$ 的直和成分生成。令 $R=\bigoplus_{m\ge0}R_m\subset\bigoplus_{m=0}^n\cO_Y(mB)$ 为由 $\bigoplus_{m=0}^n\cO_Y(mB)$ 生成的子代数。令 $Z=\Proj_Y R$，并用 $g:Z\to Y$ 表示射影。由构成方式，
\[
\Spec\hat\cO_{0,Y}\times_Y Z=\Proj_Y\Bigl(\bigoplus_{m\ge0}\hat\cO_{0,Y}(mB|_{\hat Y})\Bigr).
\]
特别地，$g$ 的例外除子不与 $g^{-1}(0)$ 相交。设 $F\subset\operatorname{Ex}(g)$ 为所有例外除子的并。此时 $0\notin g(F)$，因此 $U:=Y\setminus g(F)$ 是 $0$ 的 Zariski 开邻域。

由引理 6.2 可知 $U$ 满足 (2)。
\end{proof}

由此可知，翻转的存在性是 Euclid 拓扑意义下的局部问题。

\begin{corollaryp}{6.7}
设 $X,Y$ 为正规簇或 $\C$ 上的复解析空间，$f:X\to Y$ 为 $(K_X+D)$-翻转收缩。$f$ 的 $(K+D)$-翻转存在的充要条件如下给出。

对任意点 $y\in Y$，适当选取（Zariski 或 Euclid 的）开邻域 $y\in U_y\subset Y$，使 $f_y:f^{-1}(U_y)\to U_y$ 的 $(K+D)$-翻转存在。
\end{corollaryp}

\begin{proof}
令 $B=f_*D$。若 $f_y:f^{-1}(U_y)\to U_y$ 的 $(K+D)$-翻转存在，则 $\bigoplus_{m\ge0}\cO_{U_y}(mB|_{U_y})$ 是有限生成的，再由命题 6.6，$\bigoplus_{m\ge0}\hat\cO_{y,Y}(mB|_{\hat Y})$ 也是有限生成的。因此由命题 6.6，$\bigoplus_{m\ge0}\cO_Y(mB)$ 在 $y$ 的邻域内是有限生成的。对每个点 $y$ 这样做，即知 $\bigoplus_{m\ge0}\cO_Y(mB)$ 是有限生成代数层。
\end{proof}

下述结果在后继各节中会反复需要。（Weil 除子沿有限态射的拉回，在命题 5.20 的证明中说明。）

\begin{propositionp}{6.8}
（[Kaw88, 3.2]）设 $Y$ 与 $Y'$ 为正规且不可约的簇，$h:Y'\to Y$ 为有限满射。$B$ 为 $Y$ 上的 Weil 除子，令 $B'=h^*B$。此时
\[
\bigoplus_{m\ge0}\cO_Y(mB)\ \text{有限生成}\iff\bigoplus_{m\ge0}\cO_{Y'}(mB')\ \text{有限生成}.
\]
\end{propositionp}

% [原文 p.238]
\begin{proof}
假定 $\bigoplus_{m\ge0}\cO_Y(mB)$ 有限生成。由引理 6.2，存在满足条件 6.2(2) 的态射 $g:Z\to Y$。设 $g':Z'\to Y'$ 为 $Y'\times_Y Z$ 的正规化。它满足条件 6.2(2)。于是由引理 6.2，$\bigoplus_{m\ge0}\cO_{Y'}(mB')$ 是有限生成的。

反之，假定 $\bigoplus_{m\ge0}\cO_{Y'}(mB')$ 有限生成。为简单起见，只处理 $h$ 可分的情形。（特征为 0 时这总是成立的。）设 $q:Y''\to Y'$ 为 $Y'\to Y$ 的 Galois 闭包，$G$ 为 $Y''/Y$ 的 Galois 群，令 $B'':=q^*B$。由已证明的方向，$\bigoplus_{m\ge0}\cO_{Y''}(mB'')$ 有限生成，并由引理 6.2 得到 $g'':Z''\to Y''$。由于 $Z''$ 只有一个，$G$ 作用于 $Z''$。令 $Z:=Z''/G$ 即可。
\end{proof}

\begin{corollaryp}{6.9}
设 $f:X\to Y$ 为 $(K+D)$-翻转收缩。$g:X'\to X$ 为有限态射，$X'\to Y'\to Y$ 为 Stein 分解，其中 $X'\to Y'$ 记为 $f'$，$Y'\to Y$ 记为 $h$。以等式 $g^*(K_X+D)=K_{X'}+D'$ 定义 Q-除子 $D'$。此时，$f$ 的 $(K+D)$-翻转存在 $\iff$ $f'$ 的 $(K+D')$-翻转存在。
\end{corollaryp}

\begin{proof}
在命题 6.8 中令 $B:=f_*(K_X+D)$、$B':=f'_*(K_{X'}+D')$ 即可。
\end{proof}

上述系在 §7.3 中把某类翻转归约为翻模的存在性时，是必要的。

\begin{definition}{6.10}
设 $X$ 为正规簇。翻模收缩（flopping contraction）是指到正规簇 $Y$ 的固有双有理态射 $f:X\to Y$，它满足 $\operatorname{Ex}(f)$ 在 $X$ 内的余维数至少为 2，且 $K_X\equiv_f 0$。

若 $D$ 为 $X$ 上的 Q-Cartier Q-除子且 $-(K_X+D)$ 为 $f$-丰富，则 $f$ 的 $(K+D)$-翻转也称为 $D$-翻模（$D$-flop）。若 $f:X\to Y$ 是射线型的，则由推论 6.4(3)，$D$-翻模不依赖于 $D$ 的选取。因此，$f^+:X^+\to Y$ 或 $\varphi:X\dashrightarrow X^+$ 有时就简称为 $f$ 的翻模（flop）。

若 $(X,0)$ 是标准的（或终端的），则称翻模 $\varphi:X\dashrightarrow X^+$ 是标准的（canonical）（或终端的（terminal）），并称为标准翻模（canonical flop）（或终端翻模（terminal flop））。

% [原文 p.239]
若 $(X,0)$ 是终端的且 $D$ 是 Q-Cartier 的，则当 $0<\varepsilon\ll1$ 时，$(X,\varepsilon D)$ 仍是终端的（推论 2.35(2)）。这样，终端翻模可以看作终端翻转。然而，即使 $(X,0)$ 是标准的且 $D$ 是 Q-Cartier 的，也可能对任何 $0<\varepsilon$ 都有 $(X,\varepsilon D)$ 不是标准的。因此，标准翻模并不是标准翻转的特殊情形。
\end{definition}

% [Tutor Note: 翻模与翻转互为对偶：翻转收缩的 $K_X$ 沿收缩为负，而翻模收缩满足 $K_X\equiv_f 0$，两者都是小收缩，区别在标准类的符号。]

\begin{lemmap}{6.11}
设 $f:X\to Y$ 为翻转收缩。$B$ 为满足 $B\sim_f-2K_X$ 的约化除子，$g:X'\to X$ 为与之对应的二重覆盖（定义 2.50）。此时以下成立。
\begin{enumerate}[label=（\arabic*）,leftmargin=3.2em]
\item 由推论 6.9 定义的 $f':X'\to Y'$ 是翻模收缩。
\item 若 $(X,(1/2)B)$ 是标准的，则 $(X',0)$ 也是标准的。
\end{enumerate}
\end{lemmap}

\begin{proof}
存在非奇异开集 $U\subset X$，使 $B|_U$ 非奇异且 $X\setminus U$ 的余维至少为 2。此时由引理 2.51，$U':=g^{-1}(U)$ 非奇异。于是 $X'$ 正规，且 $g$ 沿 $B$ 分歧。由此，$K_{X'}=g^*(K_X+(1/2)B)$ 是 $f'$-平凡的。

第二个主张由命题 5.20(3) 推出。
\end{proof}

在 §6.4 与 §7.4 中，为构成某类翻转要用到 MMP。这一手法基于下述结果，它是推论 3.53 的特殊情形。

\begin{proposition}{6.12}
设 $f:X\to Y$ 为 $(K+D)$-翻转收缩。$g:X'\to X$ 为固有双有理态射，$(X',D')$ 为 $K_{X'}+D'$ 是 Q-Cartier 的正规簇对，且假定 $K_{X'}+D'-g^*(K_X+D)$ 是有效且 $g$-例外的。此时 $Y$ 上 $(X',D')$ 的拟标准模型（定义 3.50）也是 $f:X\to Y$ 的 $(K+D)$-翻转。
\end{proposition}

\section{终端翻模}
\takeaway{本节在三维情形证明终端翻模的存在性：只具有终端奇点的三维簇的翻模收缩必存在翻模（定理 6.14）；证明把问题归约到 $(\C^4,0)$ 的超曲面奇点 $u^2=f(x,y,z)$，并用到三维终端奇点的分类（推论 5.39）。}

本节的目的，是在三维情形证明终端翻模的存在性，以及极标准翻转的终止性。

\begin{remark}{6.13}
本节的讨论，在代数的设定与解析的设定下同样有效。

% [原文 p.240]
（从完全证明的观点出发，作下述指出。本节的证明基于三维终端奇点的分类（推论 5.39），而其证明用到定理 5.22。本书只对射影簇的情形给出了定理 5.22 的证明。在这个意义上，本节的证明只对三维射影簇的情形才是完全的。）
\end{remark}

% [Tutor Note: 原书“极标准翻转”对应英文底本 extremal canonical flip，“极”取“极端的”之义，与极射线同源，并非“终端（terminal）”之误，不要改写为“终端标准翻转”。]

\begin{theoremp}[三维终端翻模的存在性]{6.14}
设 $f:X\to Y$ 为只具有终端奇点的三维簇 $X$ 的 $(K_X+D)$-翻模收缩。此时 $(K_X+D)$-翻模存在。
\end{theoremp}

\begin{proof}
首先证明 $(Y,0)$ 是终端的。若 $0<\varepsilon\ll1$，则 $(X,\varepsilon D)$ 是终端的（定义 6.10）。$\overline{NE}(X/Y)$ 是 $(K_X+\varepsilon D)$-负的极锥面，故可把 $f$ 看作 $\overline{NE}(X/Y)$ 的收缩。此时由定理 3.25(4)，$K_Y$ 是 Q-Cartier 的，且 $K_X=f^*K_Y$ 成立。于是 $(Y,0)$ 是终端的（引理 2.30）。

设 $Q\in Y$ 为使 $f^{-1}(Q)$ 不是一点的点。把 $Y$ 作为解析芽 $(Q\in Y)$ 处理。由推论 6.7，只需证明 $R(Y,K_Y+f_*D)$ 是有限生成 $\cO_{Q,Y}$-代数。

由推论 5.39，$(Q\in Y)$ 具有有限覆盖 $\pi:(\tilde Q\in\tilde Y)\to(Q\in Y)$。其中 $(\tilde Q\in\tilde Y)$ 是 $(\C^4,0)$ 的超曲面奇点 $u^2=f(x,y,z)$。$\tilde Y$ 还通过态射 $(x,y,z,u)\mapsto(x,y,z)$ 成为二重覆盖 $g:(\tilde Q\in\tilde Y)\to(0\in C^3)$。设 $\iota:\tilde Y\to\tilde Y$ 为 2 阶自同构 $(x,y,z,u)\mapsto(x,y,z,-u)$。对 $\tilde Y$ 上的 Weil 除子 $F$，容易看出 $F+\iota^*F=g^*g_*F\sim0$，即 $\iota^*F\sim-F$。因此 $\iota^*R(\tilde Y,F)\cong R(\tilde Y,-F)$。

由于 $-K_X-D$ 是 $f$-丰富的，$R(Y,-K_Y-f_*D)$ 是有限生成的。由命题 6.8，$R(\tilde Y,-\pi^*f_*D)$ 是有限生成的。特别地，$R(\tilde Y,\pi^*f_*D)\cong\iota^*R(\tilde Y,-\pi^*f_*D)$ 也是有限生成的。再由命题 6.8，$R(Y,K_Y+f_*D)$ 成为有限生成。
\end{proof}

% [Tutor Note: 定理 6.14 的证明只对 $K_X$ 为 Cartier 的情形完整；一般情形要考察翻模收缩的指数 1 覆盖上的循环群作用，这又用到三维终端奇点的分类。]

上述讨论，在 $K_X$ 是 Cartier 的情形证明了下面的定理。若 $K_X$ 不是 Cartier 的，则通过考察翻模收缩的指数 1 覆盖上的循环群作用来证明。
